\chapter{Backend API Specification \& Logic}
\label{chap:backend}

The frontend is coded to expect a specific "contract" (JSON Structure) from the backend. The backend implementation must adhere to these requirements to prevent frontend crashes.

\section{Service Endpoints}

The \texttt{src/api.js} file communicates with the following endpoints:

\subsection{Route Calculation}

\textbf{Function:} \texttt{calculateRoute(start, end, travelMode)}\\
\textbf{HTTP Method:} POST\\
\textbf{Route:} \texttt{/route}

\textbf{Request Payload:}
\begin{tcolorbox}[colback=gray!10,colframe=black!50,arc=2mm]
\begin{verbatim}
{
  "start": {
    "lat": 10.8231,
    "lon": 106.6297
  },
  "end": {
    "lat": 10.7621,
    "lon": 106.6869
  },
  "travelMode": "car"
}
\end{verbatim}
\end{tcolorbox}

\textbf{Expected Response:}
\begin{tcolorbox}[colback=gray!10,colframe=black!50,arc=2mm]
\begin{verbatim}
{
  "coords": [
    {"lat": 10.8231, "lon": 106.6297},
    {"lat": 10.8200, "lon": 106.6350},
    {"lat": 10.7621, "lon": 106.6869}
  ],
  "distance_km": 12.5,
  "duration_min": 25.3
}
\end{verbatim}
\end{tcolorbox}

\textbf{Critical Fields:}
\begin{itemize}
    \item \texttt{coords}: Array of coordinate pairs for polyline rendering
    \item \texttt{distance\_km}: Total distance in kilometers
    \item \texttt{duration\_min}: Estimated travel time in minutes
\end{itemize}

\subsection{Location Search}

\textbf{Function:} \texttt{searchLocation(address)}\\
\textbf{HTTP Method:} POST\\
\textbf{Route:} \texttt{/search}

\textbf{Request Payload:}
\begin{tcolorbox}[colback=gray!10,colframe=black!50,arc=2mm]
\begin{verbatim}
{
  "address": "Ben Thanh Market"
}
\end{verbatim}
\end{tcolorbox}

\textbf{Expected Response:}
\begin{tcolorbox}[colback=gray!10,colframe=black!50,arc=2mm]
\begin{verbatim}
[
  {
    "address": {
      "freeformAddress": "Ben Thanh Market, District 1, HCMC"
    },
    "position": {
      "lat": 10.7721,
      "lon": 106.6983
    },
    "poi": {
      "name": "Ben Thanh Market"
    }
  }
]
\end{verbatim}
\end{tcolorbox}

\subsection{Camera Images}

\textbf{Function:} \texttt{fetchCameraImages(cameraId)}\\
\textbf{HTTP Method:} GET\\
\textbf{Route:} \texttt{/api/camera/\{cameraId\}/images}

\textbf{Expected Response:}
\begin{tcolorbox}[colback=gray!10,colframe=black!50,arc=2mm]
\begin{verbatim}
{
  "camera_id": "5deb576d1dc17d7c5515acf8",
  "count": 1,
  "images": [
    {
      "url": "/api/camera/5deb576d1dc17d7c5515acf8/proxy?t=...",
      "timestamp": 1763645800465
    }
  ]
}
\end{verbatim}
\end{tcolorbox}

\subsection{User Authentication}

\textbf{Function:} \texttt{login(username, password)}\\
\textbf{HTTP Method:} POST\\
\textbf{Route:} \texttt{/api/auth/login}

\textbf{Request Payload:}
\begin{tcolorbox}[colback=gray!10,colframe=black!50,arc=2mm]
\begin{verbatim}
{
  "username": "user@example.com",
  "password": "password123"
}
\end{verbatim}
\end{tcolorbox}

\textbf{Expected Response:}
\begin{tcolorbox}[colback=gray!10,colframe=black!50,arc=2mm]
\begin{verbatim}
{
  "token": "eyJhbGciOiJIUzI1NiIsInR5cCI6IkpXVCJ9...",
  "username": "user@example.com"
}
\end{verbatim}
\end{tcolorbox}

\subsection{Save Location}

\textbf{Function:} \texttt{saveLocation(name, lat, lng)}\\
\textbf{HTTP Method:} POST\\
\textbf{Route:} \texttt{/api/user/locations}

\textbf{Request Payload:}
\begin{tcolorbox}[colback=gray!10,colframe=black!50,arc=2mm]
\begin{verbatim}
{
  "name": "Ben Thanh Market",
  "lat": 10.7721,
  "lng": 106.6983
}
\end{verbatim}
\end{tcolorbox}

\subsection{Get Saved Locations}

\textbf{Function:} \texttt{getSavedLocations()}\\
\textbf{HTTP Method:} GET\\
\textbf{Route:} \texttt{/api/user/locations}

\textbf{Expected Response:}
\begin{tcolorbox}[colback=gray!10,colframe=black!50,arc=2mm]
\begin{verbatim}
[
  {
    "id": 1,
    "name": "Ben Thanh Market",
    "lat": 10.7721,
    "lng": 106.6983,
    "timestamp": "2025-11-21 10:30:00"
  }
]
\end{verbatim}
\end{tcolorbox}

\section{Algorithmic Requirements}

The backend is not just a database retrieval tool; it is a computational engine.

\subsection{Multi-Modal Routing Algorithm}

The routing engine must support three transport modes:
\begin{enumerate}
    \item \textbf{Car}: Uses road network with traffic consideration
    \item \textbf{Bike}: Uses bike lanes and roads
    \item \textbf{Walking}: Uses pedestrian paths and roads
\end{enumerate}

\textbf{Graph Construction:}
\begin{itemize}
    \item Nodes: Intersections, landmarks, transit stops
    \item Edges: Roads, bike lanes, pedestrian paths
    \item Edge Weights: Distance + Time + Safety factors
\end{itemize}

\textbf{Scoring Function:}

The backend returns a score for each route. The formula is:

\[
S = (w_1 \cdot T_{norm}) + (w_2 \cdot D_{norm}) + (w_3 \cdot S_{norm})
\]

Where:
\begin{itemize}
    \item $T$ = Travel Time (normalized)
    \item $D$ = Distance (normalized)
    \item $S$ = Safety Score (normalized)
    \item $w_1, w_2, w_3$ = Weights (configurable)
\end{itemize}

\subsection{Safety Analysis}

The system integrates traffic camera data to assess route safety:
\begin{itemize}
    \item \textbf{Camera Coverage}: Routes passing through camera-monitored areas score higher
    \item \textbf{Congestion Detection}: Real-time traffic data from cameras
    \item \textbf{Incident Reporting}: User-reported incidents on routes
\end{itemize}
